\section{Architecture}\label{sec:architecture}

\subsection{On-disk superblock (format v3)}

The superblock occupies block~0 of the device, fits in one
4096-byte block, and is the anchor of the metadata trust chain.
\Cref{fig:sb-layout} renders its v3 layout. The header carries
identification (magic, version, block size) and the global
counters (block and inode counts, free counts, table pointers).
The CRC32 covers two non-contiguous regions: \emph{region A}
(bytes 0--63, fixed-position counters and version) and
\emph{region B} (bytes 68--1688, UUID, label, RS journal, bitmap
pointer, feature flags and data protection scheme). The
\texttt{s\_crc32} field itself in bytes 64--67 is excluded from
the CRC range, as is conventional.

\begin{figure*}[t]
\centering
% Fig.1 -- FTRFS superblock layout v3 (4096 bytes)
% Standalone TikZ for inclusion via \input{}
\begin{tikzpicture}[
  every node/.style={font=\small},
  block/.style={draw, minimum height=0.8cm, anchor=west},
  scale=1
]
% --- Total block: 4096 bytes, drawn as a 14cm horizontal bar
% Scale: 1 byte = 14cm / 4096 -- but we'll group by region

% Region positions (relative widths chosen for legibility, not strict scale)
\def\xstart{0}
\def\hbar{1.0}

% Region A: bytes 0-63 (header + counters), 64 bytes
\node[block, fill=blue!20, minimum width=1.2cm] (rA) at (\xstart,0) {};
\node[font=\scriptsize] at (rA.center) {hdr};

% s_crc32: bytes 64-67, 4 bytes (excluded from CRC)
\node[block, fill=red!30, minimum width=0.4cm] (crc) at (rA.east) {};
\node[font=\scriptsize, rotate=90] at (crc.center) {crc};

% UUID + label: bytes 68-115, 48 bytes
\node[block, fill=green!20, minimum width=0.7cm] (uuid) at (crc.east) {};
\node[font=\scriptsize] at (uuid.center) {uuid};

% RS journal: bytes 116-1651, 1536 bytes (64 entries x 24 bytes)
\node[block, fill=orange!30, minimum width=4.2cm] (journal) at (uuid.east) {};
\node[font=\scriptsize] at (journal.center) {RS event journal (64$\times$24 B)};

% Misc fields: bytes 1652-1688, 36 bytes (head, bitmap_blk, features, scheme)
\node[block, fill=yellow!30, minimum width=0.7cm] (misc) at (journal.east) {};
\node[font=\scriptsize, rotate=90] at (misc.center) {feat};

% s_pad: bytes 1689-3967 (offset 1689 to before parity), 2279 bytes
\node[block, fill=gray!15, minimum width=5.4cm] (pad) at (misc.east) {};
\node[font=\scriptsize] at (pad.center) {\texttt{s\_pad[2407]} (reserved)};

% RS parity: bytes 3968-4095, 128 bytes (8 subblocks x 16 parity bytes)
\node[block, fill=purple!30, minimum width=1.4cm] (parity) at (pad.east) {};
\node[font=\scriptsize] at (parity.center) {RS parity};

% Total bracket on top
\draw[decorate, decoration={brace, amplitude=4pt}]
  ([yshift=2pt]rA.north west) -- ([yshift=2pt]parity.north east)
  node[midway, above, yshift=4pt]
  {Superblock = 4096 bytes (block 0)};

% Byte offset annotations below
\node[font=\scriptsize, below=1pt of rA, anchor=north west] {0};
\node[font=\scriptsize, below=1pt of crc, anchor=north west] {64};
\node[font=\scriptsize, below=1pt of uuid, anchor=north west] {68};
\node[font=\scriptsize, below=1pt of journal, anchor=north west] {116};
\node[font=\scriptsize, below=1pt of misc, anchor=north west] {1652};
\node[font=\scriptsize, below=1pt of pad, anchor=north west] {1689};
\node[font=\scriptsize, below=1pt of parity, anchor=north west] {3968};
\node[font=\scriptsize, below=1pt of parity, anchor=north east] {4096};

% CRC32 coverage bracket below the byte offsets
\draw[decorate, decoration={brace, mirror, amplitude=4pt}, blue!50!black]
  ([yshift=-18pt]rA.south west) -- ([yshift=-18pt]misc.south east)
  node[midway, below, yshift=-4pt, blue!50!black, font=\footnotesize]
  {CRC32 coverage: regions A (0--63) + B (68--1688), 1685 logical bytes};

% RS staging bracket
\draw[decorate, decoration={brace, mirror, amplitude=4pt}, purple!50!black]
  ([yshift=-44pt]rA.south west) -- ([yshift=-44pt]parity.south east)
  node[midway, below, yshift=-4pt, purple!50!black, font=\footnotesize]
  {RS(255,239) shortened: 8 subblocks of 211 data bytes, 16 parity bytes each (128 B parity at offset 3968)};
\end{tikzpicture}

\caption{FTRFS superblock layout, format v3 (4096 bytes).
Region A and region B are concatenated into a 1688-byte staging
buffer (1685 logical bytes plus 3 bytes of zero pad to round up
to 8 shortened RS subblocks of 211 data bytes each). The 128 bytes
of RS parity reside at offset 3968 (the trailing 128 bytes of the
block). The \texttt{s\_pad[2407]} region between the structured
header and the parity zone is reserved for future feature
extensions, in line with axis 3.1 (single source of truth: future
fields encode their offset via \texttt{offsetof()} guarded by
\texttt{BUILD\_BUG\_ON}).}
\label{fig:sb-layout}
\end{figure*}

\subsection{Universal inode protection (item 1)}

Every inode occupies 256~bytes on-disk. The first 240 bytes carry
the canonical fields (mode, uids, sizes, timestamps, flags, CRC32,
direct/indirect block pointers); the trailing 16~bytes
(\texttt{i\_reserved[0..15]}) hold the RS(255,239) parity over the
preceding 172 bytes of structured fields plus padding. The data
protection scheme \texttt{INODE\_UNIVERSAL} (value 5, item~1)
mandates this layout for every inode unconditionally; the legacy
opt-in flag \texttt{FTRFS\_INODE\_FL\_RS\_ENABLED} is preserved as
a bit definition for backward compatibility with v0.1.0/v0.2.0
images, but new images do not set it (per axis~3.3).

\subsection{On-disk allocation bitmap}

The block allocation bitmap occupies a single 4096-byte block,
formatted as 16 RS(255,239)-shortened subblocks of 239~data bytes
plus 16~parity bytes each. The total addressable bitmap covers
$16 \times 239 \times 8 = 30{,}592$ blocks, well within the bounds
of typical embedded deployments.

\subsection{Radiation event journal}

The superblock embeds a 64-entry ring buffer of
\texttt{struct ftrfs\_rs\_event} records, occupying 1536~bytes
($64 \times 24$). Each record captures the corrected block number,
a nanosecond timestamp, a symbol count, and a CRC32 over the
preceding 20~bytes. The next-write index is held in
\texttt{s\_rs\_journal\_head}. On every successful RS correction
event, the kernel writes one record (per axis~3.2). Format~v4 will
extend each record to 40~bytes by adding a Shannon entropy field
and structural sentinels (see \cref{sec:future}).

\subsection{Reed--Solomon code parameters}

FTRFS uses RS(255,239) shortened with data length tailored to the
protected region: 211 data bytes for the superblock subblocks,
239 data bytes for inode and bitmap subblocks. Each subblock
carries 16 parity bytes, tolerating up to 8 symbol errors
($\lfloor 16/2 \rfloor$) independently. Encoding and decoding are
delegated to \texttt{lib/reed\_solomon} from the Linux kernel
tree~\cite{kernelrslib}, avoiding the maintenance burden and
correctness risk of an in-tree reimplementation.
